\chapter{Macular Edema}
\index{Macular Edema}
\label{sec:Macular Edema}

\section{Overview}
\begin{itemize}[noitemsep]
  \item Macular edema is the accumulation of fluid in the central retina (macula), causing swelling and blurred central vision
  \item Diabetic macular edema (DME) is the most common cause and a leading cause of vision loss in working-age adults
  \item Cystoid macular edema (CME) often follows eye surgery, retinal vein occlusion, or inflammation
\end{itemize}
\section{Causes}
This condition happens because of breakdown of the blood--retinal barrier, allowing fluid and protein to leak from retinal capillaries into the macula. It is caused by diabetes (diabetic retinopathy), retinal vein occlusion, uveitis, cataract surgery (Irvine--Gass syndrome), age-related macular degeneration, and some medications such as prostaglandin eye drops or niacin.
\section{Risk Factors}
\begin{itemize}[noitemsep]
  \item Diabetes with poorly controlled blood sugar and hypertension
  \item Retinal vein occlusion
  \item Intraocular inflammation (uveitis)
  \item Recent eye surgery, especially complicated cataract surgery
  \item High myopia
  \item Medications that cause fluid retention in the retina
\end{itemize}
\section{Signs and Symptoms}
\subsection{Common symptoms}
\begin{itemize}[noitemsep]
  \item Blurred or wavy central vision
  \item Colors may appear washed out or distorted
  \item Difficulty reading and seeing fine detail
  \item Fluctuating vision, often worse in the morning
  \item Peripheral vision is usually preserved
\end{itemize}
\subsection{Emergency warning signs}
\begin{itemize}[noitemsep]
  \item Sudden, severe, or worsening central vision loss requires urgent eye evaluation
\end{itemize}
\section{Progression and Stages}
Macular edema may be focal or diffuse. Untreated DME tends to persist and progress, with fluid accumulation causing retinal thickening and eventual photoreceptor damage and permanent central vision loss. With early treatment, many people recover much of their vision, but chronic edema can lead to scarring and irreversible loss.
\section{Complications}
\begin{itemize}[noitemsep]
  \item Permanent central vision loss from photoreceptor damage
  \item Macular fibrosis or scarring
  \item Formation of cysts in the retina
  \item Reduced quality of life and impaired daily function
  \item Higher risk in diabetic patients of coexisting eye disease
\end{itemize}
\section{Diagnosis}
\begin{itemize}[noitemsep]
  \item Dilated fundus examination
  \item Optical coherence tomography (OCT) measures retinal thickness and fluid
  \item Fluorescein angiography or OCT angiography to identify leaking vessels
  \item Visual acuity testing
  \item Blood tests for glucose (HbA1c) in diabetic patients
\end{itemize}
\section{Prevention}
\begin{itemize}[noitemsep]
  \item Strict control of blood sugar, blood pressure, and cholesterol in diabetes
  \item Routine dilated eye screening for diabetics
  \item Timely treatment of retinal vein occlusion and inflammation
  \item Regular eye examinations after eye surgery
  \item Careful selection of medications under medical supervision
\end{itemize}
\section{Treatment Options}
\begin{itemize}[noitemsep]
  \item Intravitreal anti-VEGF injections (ranibizumab, aflibercept, bevacizumab) are first-line therapy for DME and vein-occlusion edema
  \item Intravitreal or periocular corticosteroids (dexamethasone, triamcinolone) for inflammatory or refractory edema
  \item Focal/grid laser photocoagulation for DME
  \item Treatment of the underlying cause, such as blood pressure control or anti-inflammatory therapy
  \item Surgery (vitrectomy) rarely for specific cases
\end{itemize}
\section{Living with It}
\begin{itemize}[noitemsep]
  \item Follow your injection and monitoring schedule faithfully
  \item Maintain good control of diabetes, blood pressure, and cholesterol
  \item Use good lighting and magnifiers for reading
  \item Report any sudden visual changes promptly
\end{itemize}
\section{Outlook}
The outlook has improved dramatically with anti-VEGF therapy. Most patients with DME or vein-occlusion edema gain vision with treatment, though repeat injections are usually needed. Outcomes are best with early detection and treatment; chronic untreated edema carries a poorer prognosis.
